\scanpage{294}
\mspage{210}

Nous allons prouver les implications supplémentaires

\[
b''\Longrightarrow b,\qquad c\Longrightarrow b,\qquad c''\overset{?}{\Longrightarrow}c,
\]

ce qui achèvera la démonstration.

L'implication $b''\Longrightarrow a$ est contenue dans 4.6.3. Pour l'implication $c\Longrightarrow b$, on note que la stabilité de \deleted{$\unsure{C}$} $M$ par petites $\unsure{\varinjlim|\underline{\lim}}$ équivaut à l'assertion que toute \inserted{petite} limite projective \deleted{$\unsure{\xi}$} de foncteurs $M\xrightarrow{\deleted{\unsure{F}}\unsure{\xi|\zeta}_i}\operatorname{Ens}$ représentables, est représentable. Or les objets de $M$ étant accessibles ($M$ étant accessible), les foncteurs représentables $\xi_i$ sont accessibles, donc aussi leur $\varprojlim\,\unsure{\xi|{}}$ (SGA 4 I 9.7). D'autre part, le foncteur $\xi$ est exact à gauche, comme $\varprojlim$ de foncteurs \inserted{repr. donc} exacts à gauche. Donc le fait que $\xi$ est représentable résulte du

\uline{Lemme} 4.8.11. Soit $M$ une catégorie accessible, \deleted{Alors les} \inserted{stable par $\varprojlim$ finies,} $M\xrightarrow{\xi}\operatorname{Ens}$ un foncteur. Pour que $\xi$ soit proreprésentable, il f. et s. qu'il soit exact à gauche et accessible. Donc \inserted{(SGA 4 I \unsure{8.12.2.1|8.12.8.1|§ 12.2.1})} si $M$ est stable par petites $\varprojlim$, $\xi$ est repr. \unsure{sss|ssi} il commute auxdites $\varprojlim$, et qu'il soit accessible.
% archive sheet number: 293
